% answers_short.tex —— 2013 年全国Ⅰ卷（理科） · 简短答案速查
% 由 tools/make_short_answers.py 自 answers.tex 的 \gksolhead/\gkanswer 提取，
% 并与 2013/2013_数学真题_全国1理_答案要点.md 逐题比对后生成，勿手改（改源 MD 后重跑）。

\begin{gknotice}
\textbf{使用说明}：本节只给最终答案便于快速核对，完整推导见解析版。答案为编写者独立推导，\textbf{非官方标准答案}。
\end{gknotice}

\gkpart{选择题}
\begin{center}\begin{tabular}{ccccccccccccc}
\toprule
题号 & 1 & 2 & 3 & 4 & 5 & 6 & 7 & 8 & 9 & 10 & 11 & 12 \\
\midrule
答案 & B & D & C & C & A & A & C & A & B & D & D & B \\
\bottomrule
\end{tabular}\end{center}

\gkpart{填空题}
\begin{itemize}[leftmargin=2.6em]
  \item[\textbf{13.}] $2$
  \item[\textbf{14.}] $a_n=(-2)^{n-1}$
  \item[\textbf{15.}] $-\dfrac{2\sqrt5}{5}$
  \item[\textbf{16.}] $16$
\end{itemize}

\gkpart{解答题}
\begin{itemize}[leftmargin=2.6em]
  \item[\textbf{17.}] (1) $PA=\dfrac{\sqrt7}{2}$；(2) $\tan\angle PBA=\dfrac{\sqrt3}{4}$
  \item[\textbf{18.}] (1) 见解析版；(2) $\dfrac{\sqrt{10}}5$
  \item[\textbf{19.}] (1) $\dfrac3{64}$；(2) $P(X=400)=\dfrac{11}{16}$，$P(X=500)=\dfrac1{16}$，$P(X=800)=\dfrac14$；$E(X)=\dfrac{2025}4=506.25$ 元
  \item[\textbf{20.}] (1) $\dfrac{x^2}{4}+\dfrac{y^2}{3}=1$；动圆取正半径时须去掉 $(-2,0)$，该方程描述轨迹的闭包；(2) 接触点处公切线 $x=0$ 对应 $|AB|=2\sqrt3$；另两条外公切线 $x\pm2\sqrt2y=-4$ 对应 $|AB|=\dfrac{18}7$。固定题面未限定公切线种类，弦长不唯一
  \item[\textbf{21.}] (1) $a=4,\ b=2,\ c=2,\ d=2$；(2) $k\in[1,e^2]$
  \item[\textbf{22.}] (1) 见解析版；(2) $\dfrac{\sqrt3}{2}$
  \item[\textbf{23.}] (1) $\rho^2-8\rho\cos\theta-10\rho\sin\theta+16=0$；(2) $\left(\sqrt2,\dfrac\pi4\right)$、$\left(2,\dfrac\pi2\right)$
  \item[\textbf{24.}] (1) $(0,2)$；(2) $\left(-1,\dfrac43\right]$
\end{itemize}
